\input{../../../stock/en-tete_v6.tex}

\newcommand{\titrechapitre}{titre}
\newcommand{\numerochapitre}{numéro}

\renewcommand{\headrulewidth}{0.4pt}
\renewcommand{\footrulewidth}{0.4pt}
\fancyfoot[L]{Notes de Raphaël JONTEF}
\fancyfoot[C]{\thepage}
\fancyfoot[R]{Enseirb-Matmeca}
\fancyhead[L]{Cours de Mathématiques}
\fancyhead[R]{\titrechapitre}

\begin{document}
\input{../../../stock/commands.tex}

% Page de titre custom
\setlength{\headheight}{15pt}
\begin{titlepage}
    \centering
    \vspace*{3cm}
    {\huge Chapitre \numerochapitre\par}
    {\Huge \bfseries\titrechapitre\par}
    \vspace{2cm}
    {\Large Notes rédigées par Raphaël JONTEF\par}
{\Large Avec l'aide de \par}
    \vspace{0.5cm}
    {\normalsize Cours enseigné par \par}
    \vspace{4cm}
    {\small Enseirb-Matmeca\par}
    {\small filière informatique\par}
    \vfill
    {\large \today\par}
\end{titlepage}

% plan du cours
\tableofcontents
\newpage


% -----------Corps du document--------------

On construit l'arbre minimal de hauteur $h$ comme dans la question précédente. On rajoute un pour la racine, et on intègre à gauche un sous-arbre minimal de hauteur $h-1$ et à droit un sous-arbre minimal de hauteur $h-2$.On a alors~:
$$F_h = \begin{cases*}
    1&si $h = 0$ (cas de l'arbre feuille) \\
    F_{h-1} + F_{h-2} + 1
\end{cases*}$$



\end{document}